\section{Data and Sample Construction}\label{sec:data}

The structural exercise rests on a universe-coverage administrative panel: every bid submitted on every BEC procurement auction over a four-year window, each with the firm's identifier and its registry-validated SME status. The raw extract carries \nBidsRaw~million bid-level observations across \nItemsRaw{} items, \nPurchaseOrders{} purchase orders, and \nPbu{} public buyer units---a \nDataMonths-month subset of the \itemTransactionsCount~million item-level transactions BEC intermediates over the platform throughput cited in Section~\ref{sec:institutional}. Group~65 (medical, dental, and hospital supplies) accounts for roughly \groupSixtyfivePctTx{}~percent of transactions over the window and is the treated group; the remaining \controlGroupCount{} product groups form the never-treated control set in the difference-in-differences benchmark of \ref{app:did}.

The structural analysis narrows to a sharper sample. Identification uses a symmetric window of \structuralWindowM{} months on each side of the \policyCutoffMonth{} cutoff (\nDataMonths{} months in total, the ``18m'' label of Table~\ref{tab:v3_window_sensitivity}); in Stata monthly-date units the window is $[\windowLeftStata, \windowRightStata]$, or \windowStart{} through \windowEnd{} inclusive. Within that window, I restrict to Preg\~ao auctions in Group~65. The restriction does two things. Preg\~ao is the descending-clock format that admits point identification of cost distributions from drop-out bids, so the structural estimation requires it; Group~65, in turn, is where the policy bound---the only product group to switch regime during the sample.

Bids are organized at the firm-auction level---one observation per (firm, item-auction) pair with a winner/loser flag. Three filters discipline the structural sample without re-shaping it. The bid filter $c_\epsilon = b / p^{\mathrm{ref}} \in (0, \filterCepsUpper]$ removes implausible outliers (under \nFilteredOutPct{}~percent of bids, almost always data-entry errors at the hundred-times-reference level). The participation filter $n \ge \filterMinBidders$ keeps only auctions where point identification of $F_c$ and the BNE counterfactual are both well-defined. The classification step attaches a type flag (SME or non-SME, from the BEC administrative registry validated against historical bidding to absorb transient status changes) and a pharma flag (CADMAT classes 6531, 6532, 6536, 6581---CMED-regulated pharmaceuticals and biologicals---are pharma-narrow; the remainder of Group~65 is non-pharma). The resulting structural sample carries \nObsAuctions{} firm-auction observations across \nDistinctAuctions{} distinct auctions, balanced symmetrically across Pre (\windowStart{}--February~2018, $n = \nPreAuctions{}$ auctions) and Post (\policyCutoffMonth{}--\windowEnd, $n = \nPostAuctions$). Stratum-level descriptive counts appear in Appendix Tables~\ref{tab:v3_s1_handoff} and~\ref{tab:v3_pharma_counts}.

Auction-level unobserved heterogeneity is material and not fully absorbed by the reference-price normalization. Two items inside the same CADMAT class can share a scale component that moves every bid in the auction together---volume differences, ANVISA registration burden, contract-specific complexity. I handle this with the \citet{krasnokutskaya2011} method-of-moments variance decomposition followed by best-linear-predictor shrinkage; details in Section~\ref{sec:identification}, with the recovered intraclass correlations and the variance panel in Appendix Table~\ref{tab:v3_uh_variance}. The cleaned residuals $c_\epsilon^{\text{clean}} = \exp(\hat e_{it})$ enter all downstream structural estimation.

Entry counts enter separately. For each auction I count the distinct firms of each type that submitted at least one bid, then average within (period $\times$ pharma) cells. Entry rates enter the BNE as observed equilibrium arrivals (Section~\ref{sec:identification}); the panel is Appendix Table~\ref{tab:v3_entry_rates}, with magnitudes interpreted in Section~\ref{sec:results}.

The analysis relies on four ancillary inputs. Reference prices come from the BEC catalog (\texttt{preco\_ref}) and are used both for normalization and for pricing out the fiscal cost of the rule in reais. The reference price for each item is set ex ante by the buyer from the BEC catalog and does not move with the regime change: item-level fixed effects in the DiD specifications and the structural normalization absorb cross-sectional variation in $p^{\mathrm{ref}}$ across items, and the BEC catalog architecture sets $p^{\mathrm{ref}}$ at the item level for the entire window. The variable that changes around the cutoff is regime admissibility, not item-level catalog pricing. CADMAT class codes come from the BEC product hierarchy. The historical SME flag is reconstructed from firm-level bidding histories over the full 2014--2019 period to protect against transient status changes. The pharma flag follows the CADMAT class list above.

The pair of furosemide 40~mg purchases mentioned in the introduction (item \furosItemId, a CMED-regulated generic antihypertensive) sits in the institutional sample. February~2018: \furosBidsFeb{} firms bid, clearing price \furosBelowRefPct{}~percent below reference. October~2018: \furosBidsOct{} SMEs bid, clearing price just above reference. Welfare arithmetic on this pair appears in Section~\ref{sec:welfare}; the structural averages are over \nPharmaPregaoAuc{} Preg\~ao pharmaceutical auctions in the \structuralWindowM-month sample (Table~\ref{tab:v3_s1_handoff}).
